% !TeX program = xelatex
\documentclass[11pt,a4paper]{article}
\usepackage[margin=23mm]{geometry}
\usepackage{fontspec}
\setmainfont{Latin Modern Roman}
\usepackage{amsmath,amssymb,mathtools}
\usepackage{xcolor,enumitem,booktabs,longtable,array,xurl,hyperref}
\definecolor{muted}{HTML}{53636C}
\hypersetup{colorlinks=true,urlcolor=blue,linkcolor=blue}
\setlength{\parindent}{0pt}
\setlength{\parskip}{5pt}
\newcommand{\R}{\mathbb R}
\newcommand{\E}{\mathbb E}
\newcommand{\N}{\mathbb N}
\newcommand{\ind}{\mathbf 1}
\newcommand{\Var}{\operatorname{Var}}
\newcommand{\Cov}{\operatorname{Cov}}
\newcommand{\supp}{\operatorname{supp}}
\DeclareMathOperator*{\argmin}{arg\,min}
\DeclareMathOperator*{\argmax}{arg\,max}
\newcommand{\entrymeta}[3]{{\sffamily\small\color{muted}\textbf{\texttt{#1}}\quad|\quad #2\par #3\par}}
\newcommand{\Problemref}[1]{\texttt{#1}}
\newcommand{\JELref}[2]{\texttt{#2} (JEL \texttt{#1})}
\begin{document}
\section*{Drug pricing and valuable innovation}
\textbf{Problem ID:} \texttt{I18-164}\quad \textbf{JEL:} \texttt{I18}\par
% BEGIN ECONOMICS STATEMENT
\label{problem:OP-0445}
{\sffamily\small\color{muted}Original subject group: Health economics\par}
\label{definitions:problem:30007727}
\textbf{Audit classification:} Background; proposed extension. Checked source identity and appropriate agenda/background support; expanded intervention, units, population, definedness and causal restrictions. Exact targets and variants are editorial.
\subsection*{Definitions and precise research target}
\textit{Editorial formalization.} This is a family of conditional research questions, not a single instantiated empirical estimand. Each instance supplies, before examining policy outcomes, a specification \(\theta\) containing the disease set, baseline population, baseline calendar date, and complete regimes \(p\in\{0,1\}\): a drug-price regulation and the dated reimbursement and patent regime used for comparison. Announcement and implementation timing, duration, financing, information and international responses are part of those regimes. The specification also supplies the project identities and initial states, health and cost accounting rules, probability laws and observational restrictions described below. These inputs are to be declared in an application; no particular regulation or identifying assumption is selected here.

Consider pharmaceutical firms developing treatments for the specified diseases in the United States over fifteen years from the specified starting date. For candidate development project \(j\), let \(A_j(p)\in\{0,1\}\) indicate a qualifying approval by the end of the horizon, \(N_j(p)\geq0\) discounted treated patient-years, \(b_j(p)\) independently assessed, exposure-weighted incremental quality-adjusted life-years per patient-year, and \(R_j(p)\geq0\) discounted real development and delivery cost. All these quantities and the expectation below are conditional on the fixed specification \(\theta\), suppressed in the notation. Define the innovation-value contrast
\[\Delta V=E\!\left[\sum_j\{\lambda A_j(1)N_j(1)b_j(1)-R_j(1)\}-\sum_j\{\lambda A_j(0)N_j(0)b_j(0)-R_j(0)\}\right],\]
where \(\lambda>0\) is a stated monetary valuation in common-date currency units per quality-adjusted life-year. Index a common finite universe of potential candidates across policies, including entry and abandonment; define the patient-benefit term as zero for unapproved candidates and retain incurred development costs. Require \(E[|A_j(p)N_j(p)b_j(p)|]+E[R_j(p)]<\infty\) for every \(j\) and \(p\). Identify or bound the effect on clinically valuable innovation rather than equate patent or approval counts with welfare. Explain how expected prices, financing, indication selection, and patient access enter firm decisions. Credible policy variation must separate regulation from changing science and disease demand. Report effects by unmet need and health-benefit distribution using the declared reporting definitions below.

Define a finite potential-project universe \(J\) before policy assignment, large enough to encode entry under either regime. Each project has a declared mapping to treatments and indications; shared indications or development inputs do not create additional copies of the same health gain or cost. The specification records initial development and approval states and the qualifying approval rule, including the treatment of pre-existing approvals, indication extensions and withdrawals. It distinguishes new-project effects from access effects for existing treatments. The approval indicator does not itself determine the duration of actual use.

For each project and regime, let \(\mu_{jp}\) be the finite nonnegative measure of discounted treated patient-years over the included patients, indications and times, and let \(h_{jp}\) be its measurable, integrable allocated incremental QALY rate. Define \(N_j(p)=\int d\mu_{jp}\) and, when \(N_j(p)>0\), \(b_j(p)=\int h_{jp}\,d\mu_{jp}/N_j(p)\); set \(b_j(p)=0\) when \(N_j(p)=0\). Thus heterogeneous benefits use the same exposure and discount weights as \(N_j(p)\). Negative incremental health effects are permitted. Set the health contribution to zero when \(A_j(p)=0\); no benefit rate need be assessed for such a project.

The specification supplies a common clinical benchmark portfolio, a policy-specific portfolio health outcome relative to that benchmark, and an allocation rule for project contributions satisfying
\[
 H(p)=\sum_{j\in J}A_j(p)\int h_{jp}\,d\mu_{jp},
\]
where \(H(p)\) is total discounted incremental QALYs within the declared accounting boundary. Substitution, complementarities and sequential treatments must be accounted for in this portfolio outcome; a common comparator alone does not imply additivity. The allocation rule is an input, not an assumption that standalone treatment effects add. Aggregate health value is unchanged by reallocations preserving \(H(p)\), but project and stratum reports can depend on that rule and must name it.

\(R_j(p)\) includes allocated R\&D, production and delivery resources, with common resource prices and discounting. A declared shared-cost rule must make their sum equal the included portfolio resource cost, counting failed-project expenditures and each shared input once. The specification states whether and how pre-policy sunk expenditures enter, which resource flows fall inside the fifteen-year horizon, and how terminal health benefits and their associated future delivery costs enter together. It also states how any terminal patient-years and benefits are represented in \(\mu_{jp}\) and \(h_{jp}\). The geographic boundary for health and resource costs, and the treatment of international transfers, must be declared separately from foreign responses affecting United States outcomes. Drug payments are excluded as transfers within the declared accounting boundary. This is a net-health-value objective, not a claim to measure comprehensive social welfare.

All expectations use the declared baseline population and probability laws. For each instance, specify the observable random object \(O\), its measurement and policy-assignment rules, and a substantive admissible model class \(\mathcal M_\theta\) fixed independently of outcomes. Models must generate the displayed quantities and their observable laws. Design-based restrictions and structural behavioral assumptions must be distinguished. If an equilibrium model is used, its information, feasible decisions and equilibrium selection or robust quantification must be supplied under the common conventions; none is imposed by this formulation. Identification means every admissible model with the same observed-data law yields the same displayed target; otherwise report its identified set. Integrability in each model does not imply finite bounds uniformly across models. An informative identification or bounding result must state its restrictions and establish a singleton or a proper restriction of the target values allowed before using the observed law; report explicitly when only an uninformative set is justified. No identification result is asserted here. Exact equations, horizons and assumptions are editorial, rather than source-given canonical conjectures.
\subsection*{Variants and assumption changes}
\begin{enumerate}
\item \textbf{Editorial scientific-frontier variant.} Supply predetermined classification rules and boundaries for unmet need and clinical novelty, including how unclassified projects are handled. Estimate net value within each resulting stratum rather than approval counts, using the declared portfolio benefit and shared-cost allocations. Specify whether a health-benefit distribution concerns projects, patients, patient-years or allocated health gains and supply its weighting rule; these reports need not be invariant to the allocation convention.
\item \textbf{Editorial access-versus-innovation variant.} Compare a declared short-run controlled regime holding a specified approval history fixed with endogenous fifteen-year R\&D effects. The controlled regime must specify the fixed history, its time interval, how it is imposed, and which access, indication, pricing and firm decisions remain endogenous; conditioning on observed approvals alone does not define this intervention. Supply its complete outcome-generating model or causal restrictions, terminal health-and-cost convention and international accounting boundary. This is a conditional variant, not a uniquely specified decomposition or an asserted identification result.
\end{enumerate}
\subsection*{References and source audit}
\textbf{1.} Pharmaceutical incentive mechanisms verified; exact fifteen-year net-value question not stated in inspected abstract. Locator: Journal of Economic Perspectives 39(2), 3–26; abstract. Full text or primary publication page inspected. URL: \url{https://www.aeaweb.org/articles?id=10.1257/jep.20251438}.
\begin{enumerate}\raggedright
\item\hypertarget{definitions:source:30007727:1}{} Craig Garthwaite. 2025. \emph{Economic Markets and Pharmaceutical Innovation}. Journal of Economic Perspectives 39(2), 3–26; abstract.
\url{https://www.aeaweb.org/articles?id=10.1257/jep.20251438}.
DOI: \url{https://doi.org/10.1257/jep.20251438}.
\end{enumerate}

\subsection*{Common conventions: Source mathematical conventions}

These conventions apply to each entry unless explicitly replaced.

\textbf{Models and identification.} A primitive model \(m\) specifies all probability laws, feasible choices, information, equilibrium and measurement rules used by its target. The admissible class \(\mathcal M\) is fixed independently of outcomes. If \(P_O(m)\) is its observable probability law and \(\tau(m)\) the target, its identified set at observable law \(P\) is
\[
 \mathcal I_\tau(P)=\{\tau(m):m\in\mathcal M,\ P_O(m)=P\}.
\]
Point identification means that this set is a singleton. An empty set signals incompatibility of data and assumptions. Coordinate bounds need not describe the joint set. Bounds are sharp only if every retained target is attainable or arbitrarily approachable by compatible models. Finite bounds require explicit restrictions on unobserved quantities; unrestricted confounding, hidden wealth, extrapolation or arbitrary equilibrium selection can yield uninformative sets.

\textbf{Causal interventions.} A potential outcome \(Y_i(p)\) is the outcome under a complete policy regime \(p\), including timing, financing, information and market spillovers. The target population and units are fixed before treatment. With interference, use the complete assignment vector or a justified exposure mapping; individual no-interference notation alone cannot identify an economy-wide intervention. A difference of mean potential outcomes does not require the cross-world joint distribution. Conditioning on post-treatment migration, survival, enrollment or employment changes the estimand and needs separate justification.

\textbf{Statistical statements.} An estimator, confidence set, forecast or test specifies a sampling law, model class and loss/coverage criterion. Uniform \((1-\alpha)\)-coverage means \(\inf_{m\in\mathcal M}\Pr_m\{\tau(m)\in C_n\}\geq1-\alpha\), or an explicitly asymptotic version. The whole parameter space trivially has coverage, so informativeness requires a stated width, risk or power objective. A forecasting improvement is relative to a named benchmark, information set and evaluation population.

\textbf{Equilibrium and optimization.} An equilibrium comprises feasible plans, optimality, beliefs where needed, budget constraints and all market/resource clearing. The primitives or a stated benchmark define each correspondence \(\mathcal E(p;m)\). Nonemptiness is assumed or proved, never inferred just from compact policy sets. A selected-equilibrium target uses a declared selection rule; a robust target quantifies over all permitted equilibria. Use suprema or infima unless attainment is established. An \(\varepsilon\)-optimal policy is feasible and within \(\varepsilon\) of the finite optimum value. Report an empty feasible set separately.

\textbf{Welfare and accounting.} Normative utility, interpersonal normalization, social weights, numeraire, discounting and the use of public receipts are primitives. Behavioral utility need not coincide with the welfare criterion. Household welfare includes ownership income and rents once; adding profits again to compensating variation generally double-counts them. A monetary equivalent is defined by an explicit transfer and price convention, with monotonicity and range conditions. Average effects, mechanism contributions, transfers and welfare losses are different targets. Infinite sums require convergence and dynamic feasible plans require terminal or no-Ponzi conditions.

\textbf{Source versions.} A working-paper date, revision date and journal publication year are distinguished. A DOI for a journal version is not automatically the DOI of an earlier manuscript. Locators refer to the stated version and printed pages where specified. A metadata-only check does not verify a claim about a full-text conclusion. Every entry's source subsection narrows the provenance claim accordingly.
% END ECONOMICS STATEMENT
\end{document}
